
\documentclass{article}
\usepackage{colortbl}
\usepackage{makecell}
\usepackage{multirow}
\usepackage{supertabular}

\begin{document}

\newcounter{utterance}

\centering \large Interaction Transcript for game `cladder', experiment `full\_v1.5\_default', episode 6482 with llama{-}3.1{-}8b{-}instruct.
\vspace{24pt}

{ \footnotesize  \setcounter{utterance}{1}
\setlength{\tabcolsep}{0pt}
\begin{supertabular}{c@{$\;$}|p{.15\linewidth}@{}p{.15\linewidth}p{.15\linewidth}p{.15\linewidth}p{.15\linewidth}p{.15\linewidth}}
    \# & \multicolumn{2}{c}{Player} && \multicolumn{2}{c}{Game Master} \\
    \hline

    \theutterance \stepcounter{utterance}  
    & & & \multicolumn{4}{p{0.6\linewidth}}{
        \cellcolor[rgb]{0.9,0.9,0.9}{
            \makecell[{{p{\linewidth}}}]{
                \texttt{\tiny{[P1$\langle$GM]}}
                \texttt{You are an expert in causal inference. The following question is not a typical commonsense query, but rather a meticulously designed question created by a professor specializing in causal inference, intended to assess the students' mastery of the course content.} \\
\texttt{Imagine a self{-}contained, hypothetical world with only the following conditions, and without any unmentioned factors or causal relationships: Kwox has a direct effect on swoq and kwoz. Swoq has a direct effect on kwoz. For those who are not kwox, the probability of kwoz is 36\%. For those who are kwox, the probability of kwoz is 56\%. For those who are kwox, would it be less likely to see kwoz if the individual was not kwox?} \\
\texttt{Start your answer with 'Yes' or 'No', followed by additional reasoning or evidence to support your explanation.} \\
            }
        }
    }
     \\ \\

    \theutterance \stepcounter{utterance}  
    & \multicolumn{4}{p{0.6\linewidth}}{
        \cellcolor[rgb]{0.561,0.737,0.561}{
            \makecell[{{p{\linewidth}}}]{
                \texttt{\tiny{[P1$\rangle$GM]}}
                \texttt{No.} \\
\\ 
\texttt{To answer this question, we need to apply the concept of counterfactuals and potential outcomes. We are comparing the probability of kwoz for individuals who are kwox to the probability of kwoz for individuals who are not kwox, but under the counterfactual scenario where the individual is not kwox.} \\
\\ 
\texttt{Let's denote the probability of kwoz for individuals who are kwox as P(kwoz|kwox) = 0.56, and the probability of kwoz for individuals who are not kwox as P(kwox) = 0.36.} \\
\\ 
\texttt{We want to compare P(kwoz|not kwox) to P(kwoz|kwox). However, we are given the probability of kwoz for individuals who are kwox, but not the probability of kwox for individuals who are not kwox. To make this comparison, we need to apply the law of total probability.} \\
\\ 
\texttt{Let's denote the probability of being kwox as P(kwox) and the probability of not being kwox as P(not kwox). We can write the law of total probability as:} \\
\\ 
\texttt{P(kwoz) = P(kwoz|kwox)P(kwox) + P(kwoz|not kwox)P(not kwox)} \\
\\ 
\texttt{We are given P(kwoz|kwox) = 0.56 and P(kwoz|not kwox) = 0.36. We also know that P(kwox) + P(not kwox) = 1.} \\
\\ 
\texttt{To find P(kwoz), we need to know the value of P(kwox). However, we are not given this value. Nevertheless, we can still answer the question by considering the relationship between P(kwoz|not kwox) and P(kwoz|kwox).} \\
\\ 
\texttt{If P(kwoz|not kwox) < P(kwoz|kwox), then it would be less likely to see kwoz if the individual was not kwox. However, we are given that P(kwoz|not kwox) = 0.36 and P(kwoz|kwox) = 0.56. Since 0.36 < 0.56, it would be less likely to see kwoz if the individual was not kwox.} \\
\\ 
\texttt{However, the question asks whether it would be less likely to see kwoz if the individual was not kwox, given that they are kwox. This is a counterfactual question, and we need to consider the potential outcomes under the counterfactual scenario.} \\
\\ 
\texttt{Let's denote the potential outcome of kwoz for individuals who are kwox as Y(kwox) and the potential outcome of kwoz for individuals who are not kwo} \\
            }
        }
    }
    & & \\ \\

    \theutterance \stepcounter{utterance}  
    & & & \multicolumn{2}{p{0.3\linewidth}}{
        \cellcolor[rgb]{0.9,0.9,0.9}{
            \makecell[{{p{\linewidth}}}]{
                \texttt{\tiny{[GM$|$GM]}}
                \texttt{No} \\
            }
        }
    }
    & & \\ \\

    \theutterance \stepcounter{utterance}  
    & & & \multicolumn{2}{p{0.3\linewidth}}{
        \cellcolor[rgb]{0.9,0.9,0.9}{
            \makecell[{{p{\linewidth}}}]{
                \texttt{\tiny{[GM$|$GM]}}
                \texttt{yes} \\
            }
        }
    }
    & & \\ \\

    \theutterance \stepcounter{utterance}  
    & & & \multicolumn{2}{p{0.3\linewidth}}{
        \cellcolor[rgb]{0.9,0.9,0.9}{
            \makecell[{{p{\linewidth}}}]{
                \texttt{\tiny{[GM$|$GM]}}
                \texttt{game\_result = LOSE} \\
            }
        }
    }
    & & \\ \\

\end{supertabular}
}

\end{document}
